% preamble.tex — LaTeX header includes for the Pandoc PDF build.
% This file is injected via include-in-header, i.e. AFTER Pandoc's own
% highlighting macros, so the redefinitions of Shaded/shadecolor below
% deliberately override the defaults.

% --- Page geometry: narrow margins (1cm sides, 2cm top/bottom). ----------
\usepackage[a4paper,top=2cm,bottom=2cm,left=2cm,right=2cm,
            includehead,headsep=8pt]{geometry}
\usepackage{microtype}

% --- Numbering: the Markdown headings already carry manual numbers (1.1,
% 4.2, ...) for the web build, so suppress LaTeX's automatic section and
% chapter numbering to avoid the doubled "2.1 1.1 ..." effect. Figures are
% then numbered continuously (1, 2, 3, ...) rather than per (unnumbered)
% chapter.
\setcounter{secnumdepth}{-1}
\usepackage{chngcntr}
\counterwithout{figure}{chapter}

% --- Figures keep their natural size (the TikZ figures are drawn at it);
% Pandoc's \pandocbounded already shrinks one wider than the text. -------
\usepackage{graphicx}
% the captions carry their hand-written number (Figure 3.1.), as the web
% build shows them, so LaTeX adds no label of its own
\usepackage{caption}
\captionsetup[figure]{labelformat=empty}
\makeatletter
\renewcommand{\fps@figure}{htbp}
% the page numbers of the table of contents take four digits past page 999
\renewcommand{\@pnumwidth}{2em}
\renewcommand{\@tocrmarg}{3em}
\makeatother

% --- Breakable URLs and inline code, so long paths/links do not run past
% the right margin. (Inline code gets \allowbreak hints from break-code.lua;
% xurl makes \url/\href break at any character.)
\usepackage{xurl}
\setlength{\emergencystretch}{3em}

% --- Code blocks: wrap long lines, and frame them on a light background so
% they stand out from the surrounding prose instead of washing into it. ---
\usepackage{fvextra}
\fvset{breaklines=true,breakanywhere=true,fontsize=\small}

% The frame is a breakable tcolorbox, which splits a block at the end of a
% page within the text height (the framed package overran it by up to 2.5pt
% when it split a block).
\usepackage{xcolor}
\usepackage[breakable,skins]{tcolorbox}
\definecolor{shadecolor}{RGB}{244,244,248}   % very light blue-gray
\definecolor{codeframe}{gray}{0.55}
% a chapter built alone may have no code block, and then no Shaded
\ifdefined\Shaded\else\newenvironment{Shaded}{}{}\fi
\renewenvironment{Shaded}{%
  \begin{tcolorbox}[breakable,enhanced jigsaw,colback=shadecolor,
    colframe=codeframe,boxrule=0.4pt,arc=0pt,outer arc=0pt,
    left=5pt,right=5pt,top=5pt,bottom=5pt,boxsep=0pt,
    before skip=6pt,after skip=6pt]}%
 {\end{tcolorbox}}

% --- A discreet running header. ------------------------------------------
\usepackage{fancyhdr}
\pagestyle{fancy}
\fancyhf{}
\fancyhead[L]{\small ADS — A User Manual}
\fancyhead[R]{\small\thepage}
\renewcommand{\headrulewidth}{0.2pt}
\setlength{\headheight}{14pt}

% --- A few mathematical/relational Unicode symbols appear in inline code
% and ASCII diagrams (e.g. "Ax <= b" written with the glyph). Map them so
% the default pdflatex engine can typeset them in text and verbatim alike.
\DeclareUnicodeCharacter{2264}{\ensuremath{\le}}      % <=
\DeclareUnicodeCharacter{2265}{\ensuremath{\ge}}      % >=
\DeclareUnicodeCharacter{221E}{\ensuremath{\infty}}   % infinity
\DeclareUnicodeCharacter{2208}{\ensuremath{\in}}      % element of
\DeclareUnicodeCharacter{2192}{\ensuremath{\rightarrow}} % ->
\DeclareUnicodeCharacter{2212}{\ensuremath{-}}        % minus sign
\DeclareUnicodeCharacter{2026}{\ldots}                % ellipsis
\DeclareUnicodeCharacter{25CF}{\textbullet}           % black circle
\DeclareUnicodeCharacter{03B2}{\ensuremath{\beta}}    % beta (an edition 0.1G-beta)
